\documentclass[10pt,a4paper]{amsart}
\usepackage[margin=19mm]{geometry}
\usepackage{amsmath,amssymb,mathtools}
\usepackage[scr=rsfs,cal=euler]{mathalfa}
\usepackage{xfrac}
\usepackage[unicode,hidelinks,bookmarks=false]{hyperref}
\newcommand{\ie}{i.e.\ }
\newcommand{\cf}{cf.\ }
\newcommand{\resp}{respectively\ }
\newcommand{\Subsection}{\subsection}
\providecommand{\nobreakdash}{\nobreak\hbox{-}}
\setlength{\emergencystretch}{2em}
\multlinegap0pt
\usepackage{amssymb}
\usepackage{mathtools}
\usepackage{xcolor}
\usepackage{tikz}
\usepackage{enumerate}
\usepackage{bm}
\newtheorem{proposition}{Proposition}
\newcommand{\RR}{\ensuremath{\mathbb{R}}}
\title{Koszulity of a certain dioperad}
\author{Alex Takeda}
\date{Source: arXiv:2511.02829v2 (CC BY 4.0)}
\begin{document}
\maketitle

\noindent\emph{Source and licence.} Koszulity of a certain dioperad. Version arXiv:2511.02829v2 (CC BY 4.0), distributed by arXiv under the Creative Commons Attribution 4.0 International licence, http://creativecommons.org/licenses/by/4.0/, as recorded in the arXiv licence metadata for this exact version and shown on the abstract page https://arxiv.org/abs/2511.02829v2. Full licence notice: \texttt{licenses/arxiv-2511.02829-CC-BY-4.0.txt}.\par\smallskip

\noindent\textit{Editorial scope.} Counted unit: the article's proposition prop:contractible (source0.tex lines 656--660). The article's complete original proof of it is delivered with it (source0.tex lines 661--675, one \texttt{proof} environment of 15 lines). No mathematics is written, repaired or re-derived: the statement and the proof are the article's own lines, in the article's own notation, and no retained line is substituted or re-flowed.  No further source-local context is needed: every cross-reference of every retained line resolves inside the two delivered spans.  Nothing was omitted from any retained span, so the package carries no omission record.  No retained line contains a citation, so the wrapper carries no bibliography and no bibliography entry.  No retained line contains a figure environment, an image-inclusion command or drawing code, so no image is delivered and the package declares no asset dependency.  Editorial formatting only, in addition: an AMS \texttt{amsart} preamble that restates the article's own declarations and macros used by the retained lines, namely the environments \texttt{proposition} on separate counters, together with the standard packages those lines need.  The retained environments print in this document as Proposition 1 in order of appearance, whereas the article prints the same items with the numbering of its own full document.  The complete native source of the article as distributed in the arXiv e-print for this exact version is bundled unchanged as \texttt{sources/188450/source0.tex}.
\begin{proposition}\label{prop:contractible}
    For any collection of pairwise distinct equivalence classes 
    \[ \mathbf{C} = \{C_1,\dots,C_r\}, \] 
    the space $\mathrm{Clov}_\mathbf{C}$ is contractible if and only if there are $r$ cuts $\gamma_1,\dots,\gamma_r$, one in each equivalence class, that are pairwise disjoint and divide the complex plane into regions such that each contains at least one of the $k$ output directions. In every other case, $\mathrm{Clov}_\mathbf{C}$ is empty.
\end{proposition}

\begin{proof}
    Let us start by proving the contrapositive of the last statement. If $\mathrm{Clov}_\mathbf{C}$ is non-empty, $\mathrm{ClovQ}_\mathbf{C}$ must contain at least one point; taking the $r$ saddle cuts gives us the desired representatives $\gamma_1,\dots,\gamma_r$.

    For the converse, let $\gamma_1,\dots,\gamma_r$ satisfy the condition in the statement. We can easily construct a regularized metric tree whose quadratic differential is cloven by $\mathbf{C}$, say, by picking a single internal vertex for each region, length $2$ for all internal edges and length $1$ for all of the outgoing $k$ leaves. This proves that $\mathrm{ClovQ}_\mathbf{C}$ is not empty.

    Let us now prove contractibility. We first note that the image of $\mathrm{ClovQ}_\mathbf{C}$ under the projection
    \[ \pi \colon Q^\mathrm{Str}_\mathrm{reg}(k;i_1,\dots,i_k) \twoheadrightarrow Q^\mathrm{Str}(k + i_1 + \dots + i_k + 2) \]
    is an open ball, being a product of the metric tree spaces $\mathrm{MetTr}(n_a)$ for some numbers $\{n_a\}$ for $a \in \{1,\dots,r+1\}$, satisfying
    \[ \sum_a n_a = k + i_1 + \dots + i_k + 2r, \]
    where each $n_a$ is the number of leaves of each subtree, once we cleave along the saddle cuts. So $\mathrm{ClovQ}_\mathbf{C}$ is a subspace of
    \[ \pi(\mathrm{ClovQ}_\mathbf{C}) \times (\RR^k)/\RR_{>0} \cong \pi(\mathrm{ClovQ}_\mathbf{C}) \times \Delta^{k-1} \]
    that is, the (trivial) bundle over an open ball with fiber the $(k-1)$-simplex. In each fiber, this subspace is cut out by a certain finite system of linear equations, each one of the form
    \[ \min_{i \in I} (\lambda_i+d_i) < \min_{j \in J} (\lambda_j + d_j) \]
    where $I$ and $J$ are subsets of the outgoing leaf indices $\{1,\dots,k\}$ corresponding to two neighboring regions among the $(k+1)$ regions cut out by the critical graph, and $d_i$ are positive reals, fixed by the metric on the internal edges. This is an intersection of open convex polyhedral subsets of $\Delta^{k-1}$, and therefore an open convex polyhedron. Moreover, this polyhedron in the fiber varies continuously along the base $\pi(\mathrm{ClovQ}_\mathbf{C})$, which proves that $\mathrm{ClovQ}_\mathbf{C}$ is homeomorphic to an open ball, and therefore the corresponding subcomplex $\mathrm{Clov}_\mathbf{C} \subset  Z_{(k;i_1,\dots,i_k)}$ is contractible.
\end{proof}

\end{document}
